\tituloPreliminar{HISTORIAL DE CAMBIOS}{HISTORIAL DE CAMBIOS}
\begingroup
\fuenteTabla\fontsize{9.3bp}{11bp}\selectfont
\setlength{\tabcolsep}{4bp}
\hyphenpenalty=10000 \exhyphenpenalty=10000 \sloppy
\renewcommand{\arraystretch}{1.2}
\begin{longtable}{|>{\centering\arraybackslash}p{0.16\linewidth}|>{\centering\arraybackslash}p{0.13\linewidth}|>{\centering\arraybackslash}p{0.18\linewidth}|p{0.23\linewidth}|>{\centering\arraybackslash}p{\dimexpr0.30\linewidth-10\tabcolsep-6\arrayrulewidth\relax}|}
\hline
\textbf{Versión} & \textbf{Fecha} & \textbf{Sección del documento modificada} & \centering\arraybackslash\textbf{Descripción de cambios (corta)} & \textbf{Responsable (S)}\\\hline
\endfirsthead
\hline
\textbf{Versión} & \textbf{Fecha} & \textbf{Sección del documento modificada} & \centering\arraybackslash\textbf{Descripción de cambios (corta)} & \textbf{Responsable (S)}\\\hline
\endhead\textcolor{red}{Se considerará como versión 1.0 la entregada al final del ciclo 2026-30 al director del proyecto de grado} & 04/11/2026 & Estructura general y formato & Creación de la estructura por capítulos y secciones; adaptación del formato y ajuste de márgenes simétricos. & \textcolor{red}{Daniel Felipe Ramírez Vargas} \\\hline
 & \textcolor{red}{19/11/2026} & {\raggedright\textcolor{red}{• Portada\newline • 2.1 Perspectiva del producto}\par} & \textcolor{red}{• Redacción de la perspectiva del producto.\newline • Datos, autores y logo de portada.} & \textcolor{red}{Daniel Felipe Ramírez Vargas} \\\hline

\caption{Historial de cambios}\label{tab:historial}
\end{longtable}
\endgroup
% Registrar cambios relevantes agrupados, no cada ajuste menor.

